\section{An intrinsic normal current and its local-unit harmonic drift}
\label{sec:r18-intrinsic-normal-current}

\paragraph{Status and contribution boundary.}
The statements here are \texttt{DIRECT\_COROLLARY} consequences of the
local operator LO, its exact energy identity, the separately normal
stratum tangency lemma RS0, and the frozen original smooth presentation
quotient SG8.  They give a genuine original symmetric order-zero current
for each fixed mixed tail of a full polar.  The current is separately
normal, with explicitly retained boundary measures.  Separate closedness
is \emph{not} asserted.  Its boundary is exactly a local-unit harmonic
defect which the printed global full-polar axioms have not yet been
shown to annihilate.  No general signed PL relation descent follows.

The duality and analytic-descent researchers independently developed
the weighted normality route and then compared its algebra; they are
construction contributors to this section.  Author-separated review is
recorded separately.  The root researcher checked its scope and signs
but did not supply a proof of the missing harmonic gate.

Keep the original \(U,F,L,C,\Lambda,b\) and the fixed arbitrary global
DC tails \(v_3,\ldots,v_b\) of LO.  Work on a pure-dimensional
component of \(U\) of dimension \(n\ge1\); the construction is local
on components.  The local scalar operation \(\mathscr L\), signed
biderivation \(\Gamma\), drift measures \(a_i\), and symmetric
Hessian coefficient measures \(B_{ij}\) are exactly LO6--LO17.

All real calculations use the two independent exterior columns of
HF9, with no cross-column sign.  In bidimension \((1,1)\) the
normalization factor in HF10 is one.  When we display test boundary
operators, we define them explicitly by
\[
       (\partial_LT)(\alpha)=T(d_L\alpha),\qquad
       (\partial_RT)(\beta)=T(d_R\beta)
 \tag{NC1}
\]
for bidegrees \((0,1)\) and \((1,0)\), respectively.  The
corresponding graded signs in an original-current differential convention
can then be applied without changing any vanishing or normality claim.

\subsection{Weighted real pushforwards have finite separate boundaries}

Let \(t:V\to\mathbb R^m\) be any local actual unit moment on an
ordinary open \(V\subset U\).  The coefficients supplied by LO8
are compatible under protected restrictions, and hence are locally
finite measures on \(C\cap V\).  Choose any real original
\(\xi\in A_c^{0,0}(V)\), with no requirement that it be a function
of \(t\).  Define the compact real matrix measure and current
\[
 \mu_{ij}^\xi=t_*(\xi B_{ij}),\qquad
 T_{t,\xi}(\eta)=\sum_{i,j}\int_{\mathbb R^m}\eta_{ij}\,d\mu_{ij}^\xi.
 \tag{NC2}
\]
Only compactly supported analytic coefficient measures are pushed, so
no ambient properness assertion is needed.

\begin{proposition}[Exact weighted divergence identity]
For every \(j\), in the distributional sense on \(\mathbb R^m\),
\[
       \sum_i\partial_i\mu_{ij}^\xi
            =t_*\bigl(\xi a_j+\Gamma(\xi,t_j)\bigr).
 \tag{NC3}
\]
The right side is a finite compactly supported signed Radon measure.
By symmetry the same identity applies to the other tensor column.
Thus \(T_{t,\xi}\) is a compact separately normal real double current
of bidimension \((1,1)\).
\end{proposition}

\begin{proof}
Fix \(\phi\in C_c^\infty(\mathbb R^m)\).  Choose an original
compact smooth \(\rho\) in \(V\), equal to one near
\(\operatorname{supp}\xi\), and put
\(u=\widetilde{\rho t_j}\),
\(h=\widetilde{\xi\phi(t)}\), extending both by zero to \(U\).
These are compact original smooth functions and hence are in
\(\mathcal D(U)\).  LO17 and the locality of \(\mathscr L\)
and \(\Gamma\) give
\[
             \int_{C\cap V}\xi\phi(t)\,da_j
                      =-\int_{C\cap V}\Gamma(\xi\phi(t),t_j).
 \tag{NC4}
\]
Indeed \(u=t_j\) near the support of \(h\); every difference
outside that neighborhood pairs to zero by the proved carrier-germ
locality.  The derivation identity LO16 and the local formula LO15
expand the right side as
\[
 -\int\phi(t)\,d\Gamma(\xi,t_j)
                   -\sum_i\int\xi\,\phi_i(t)\,dB_{ij}.
 \tag{NC5}
\]
Rearranging (NC4)--(NC5) is exactly (NC3), with the usual negative
sign in the definition of a distributional derivative.
The measures \(\xi a_j\) and \(\Gamma(\xi,t_j)\) have compact
support in \(C\cap\operatorname{supp}\xi\), by locality, and
are finite there.  Their pushforwards are therefore finite.  Since
\(B_{ij}=B_{ji}\), the two column divergences have the same type
of bound.  With the test-boundary convention (NC1), the real
boundary coefficients are the negatives of the measures in (NC3).
No cutoff derivative has been omitted: it is the entire explicit
term \(\Gamma(\xi,t_j)\).
\end{proof}

For an actual affine refinement \(t=At'+c\) on \(V\), LO10 and
ordinary coefficient pushforward yield
\[
                        p_*^{\rm tensor}T_{t',\xi}=T_{t,\xi},
             \qquad p(y')=Ay'+c,
 \tag{NC6}
\]
with no rank hypothesis.  This comparison uses the same actual
analytic coefficient \(\xi\); it does not transport a full polar
between ordinary domains.

\subsection{The whole original smooth form quotient is respected}

\begin{lemma}[Zero smooth presentations are annihilated with arbitrary analytic weights]
Let \((t,\eta)\) be a smooth \((1,1)\) presentation on \(V\),
and let \(\xi\in A_c^{0,0}(V)\).  If its original smooth form
is zero on an ordinary neighborhood of
\(C\cap\operatorname{supp}\xi\), then
\[
                         T_{t,\xi}(\eta)=0.
 \tag{NC7}
\]
Real compact cutoffs equal to one near the compact coefficient
support may be inserted in \(\eta\).
\end{lemma}

\begin{proof}
Cover \(C\cap\operatorname{supp}\xi\) by finitely many original
compact affinoid domains \(K_\ell\) contained in the given zero
neighborhood and in \(V\).  The exact full-image smooth quotient
SG8 and its sheafification argument imply that \(\eta\) restricts
to zero on each full image \(t(K_\ell)\).  These images are compact
finite polyhedral sets by the same original moment-image theorem
used in SG.  A common finite subdivision makes their union
\(Z\) a finite polyhedral set with
\[
            \operatorname{supp}T_{t,\xi}
                      \subset t(C\cap\operatorname{supp}\xi)\subset Z.
\]
By (NC3), the real current has finite separate boundary measures.
RS0 therefore makes its coefficient density tangential in both
columns on every relative open cell of the chosen subdivision.
The restriction of \(\eta\) to each such cell is zero, so its
pairing with the corresponding Borel restriction of the coefficient
measures is zero.  Finite addition proves (NC7).  The argument retains
all mass on lower-dimensional faces; it neither discards face terms
nor declares the restrictions themselves closed.  All forms paired
here are smooth, so no product of two singular currents has been
introduced.
\end{proof}

\begin{theorem}[An intrinsic symmetric separately normal current]
There is a unique original symmetric order-zero current
\(\mathcal T_{\Lambda;v_3,\ldots,v_b}\) of bidimension \((1,1)\)
on \(U\), supported on \(C\), with the local presentation rule
\[
 \mathcal T\bigl(\xi\,t^*\eta\bigr)
                         =T_{t,\xi}(\eta)
                         =\sum_{i,j}\int_{C\cap V}
                                                \xi\,\eta_{ij}(t)\,dB_{ij}
 \tag{NC8}
\]
for every original compact smooth scalar \(\xi\) in a local
actual moment chart and every smooth real \((1,1)\)-form \(\eta\).
It is defined on the full original compact smooth CLD test sheaf,
with every actual local moment refinement.  It satisfies
\[
       \mathcal T(d'h\otimes d''u)=\int\Gamma(h,u)
 \tag{NC9}
\]
whenever the displayed form has compact support in a protected local
chart, with the same formula after finite ordinary partition of unity.

For compact original tests of the form
\(\alpha=\sum_j h_jd''t_j\) and
\(\beta=\sum_i k_i d't_i\), its test boundaries are
\[
          (\partial_L\mathcal T)(\alpha)=-\sum_j\int h_j\,da_j,
          \qquad
          (\partial_R\mathcal T)(\beta)=-\sum_i\int k_i\,da_i.
 \tag{NC10}
\]
They are original coefficient-Radon currents.  Thus \(\mathcal T\)
is separately normal; the boundary measures have not been set to zero.
\end{theorem}

\begin{proof}
Consider first two local presentations of the same original smooth
form.  On their overlap, pass to the actual joint unit moment;
(NC6) identifies the projected coefficient pairings.  Their difference
is a zero smooth presentation, annihilated by (NC7) with each arbitrary
compact analytic weight.  Thus the local rule (NC8) is independent
of the moment presentation.  If an auxiliary scalar cutoff equal to
one near the compact support of an actual form is changed, its
difference multiplies a zero actual form; (NC7) again gives independence.
The rule consequently glues through a finite original smooth partition
of unity on each compact test support.  This is the ordinary sheaf
construction for currents, with compatibility proved before gluing.

On every fixed local test stage, the finite variations of \(B_{ij}\)
bound the pairing by finitely many coefficient supremum norms.  This
proves order-zero continuity in the original test topology.  The
coefficient measures are supported on \(C\), giving the asserted
support.  Symmetry is \(B_{ij}=B_{ji}\) in the HF two-column
convention.  Formula (NC9) is LO15 in a common actual moment, and
then follows globally by the proved presentation compatibility.
All charts in this construction are original local charts; none of
their units is assumed to extend to the fixed larger \(U\).

For one compact coefficient \(h\) in a local chart, repeat (NC4)
with \(\xi=h\) and the real test equal to one near the relevant
compact image.  Equivalently apply LO17 with a compact extension of
\(t_j\) equal to it near \(\operatorname{supp}h\).  This gives
\[
                    \int\Gamma(h,t_j)=-\int h\,da_j.
 \tag{NC11}
\]
The differential of \(\alpha=\sum_j h_jd''t_j\) is
\(\sum_j d'h_j\otimes d''t_j\) in the specified convention.
Equations (NC9) and (NC11) prove the first formula (NC10).  The
second follows in the same way from symmetry.  Arbitrary compact
original tests of these bidegrees admit finite such local
presentations and a subordinate partition of unity, so the formulas
prove coefficient-Radon bounds for both boundaries on every fixed
original test stage.  Their well-definedness on the original quotient
also follows directly because they are boundaries of the already
defined original current \(\mathcal T\).  This completes the
normality assertion without presupposing it on the analytic space.
\end{proof}

\subsection{The exact unresolved closedness condition}

For this fixed full polar and tail, the following are equivalent:
\[
 \begin{gathered}
 \partial_L\mathcal T=0\quad\hbox{and}\quad\partial_R\mathcal T=0;\\
 \mathscr L(\log|g|)=0\quad\hbox{on every ordinary local domain
 on which }g\hbox{ is an original analytic unit}.
 \end{gathered}
 \tag{NC12}
\]
Indeed vanishing of all local drift measures makes (NC10) zero on
every test.  Conversely, one zero boundary tested against
\(h\,d''\log|g|\), with all compact original smooth \(h\), kills
the corresponding drift measure by Radon uniqueness.  Symmetry then
kills the other boundary as well.

The full global harmonic generator proves this only when the local
unit logarithm agrees along the local carrier with a globally
\(F\)-harmonic element of \(\mathcal D(U)\); whole-domain unit
moments are one such case.  The general local-unit condition (NC12)
is \texttt{OPEN} for arbitrary full polars in the present argument.
No counterexample satisfying the entire original full-polar ideal has
been supplied.  In particular the weighted-cutoff counterexample FC6
is not a failure of (NC12): it concerns a different, nonsymmetric
ordered functional.

Separate normality does not justify setting the double boundary to
zero.  The scalar identity \(\int\mathscr L(h)=0\) and the
energy identity merely give the usual cancellation between the drift
term and the divergence of the second-order coefficient; they do not
supply an additional independent equation.  No use of
\(\partial_L\partial_R\mathcal T=0\) is made here.

\paragraph{Why this remains short of the signed analytic relation theorem.}
The construction solves a smooth localization interface: a fixed
original full polar determines an actual order-zero current with
finite separate boundaries, on every original local smooth chart and
actual refinement.  It neither identifies that current as an allowed
positive background nor proves that its drift vanishes.  Even
separate closedness, if later proved, would not by itself identify a
zero curvature relation with a zero primitive pairing.  Singular
PL boundary cancellation and the appropriate primitive or exactness
bridge would still require proof.  No convolution is declared
zero-preserving, no norm level replaces the absolute carrier, and
no ordinary-current identity is promoted to all GK delta tests.
KC12 and variable-PS/CTO remain \texttt{OPEN}.

\paragraph{Exact dependencies.}
LO6--LO17 supply the original analytic coefficient measures, all local
actual affine covariance, the biderivation and compact energy identity.
RS0 is the explicit tangency lemma for finite separate boundaries.
SG8--SG9 supply only the full-image smooth presentation definition,
compact original moment-image polyhedrality and its ordinary
sheafification argument.  The present proof uses none of the
unproved signed analytic relation quotient or full-polar localization
claims.  Exact source hashes and author-separated reviews are recorded
in the adjacent ledger.
